% Generated draft: requires mathematical and source review.
\documentclass[11pt]{article}
\usepackage[margin=25mm]{geometry}
\usepackage{fontspec}
\setmainfont{Latin Modern Roman}
\usepackage{amsmath,amssymb,mathtools,unicode-math}
\setmathfont{Latin Modern Math}
\usepackage{enumitem,needspace,xurl,hyperref}
\setlength{\parindent}{0pt}
\setlength{\parskip}{5pt}
\setlength{\emergencystretch}{3em}
\newcommand{\N}{\mathbb{N}}
\newcommand{\Z}{\mathbb{Z}}
\newcommand{\Q}{\mathbb{Q}}
\newcommand{\R}{\mathbb{R}}
\newcommand{\C}{\mathbb{C}}
\newcommand{\F}{\mathbb{F}}
\newcommand{\A}{\mathbb{A}}
\newcommand{\PP}{\mathbb{P}}
\newcommand{\E}{\mathbb{E}}
\newcommand{\eps}{\varepsilon}
\newcommand{\dd}{\,\mathrm{d}}
\begin{document}
\section*{UnsolvedMath ID: 30003940}
\textbf{Source problem code:} catalogue entry
\subsection{Definitions and mathematical statement}
The supplied record concerns a two-species Keller--Segel chemotaxis model in two spatial dimensions. It denotes one model parameter by \(\chi_2\) and asks about the qualitative regime
\[
\chi_2\ll 1.
\]
The equations, the definition and normalization of \(\chi_2\), its admissible values, and the parameters held fixed in this regime are not supplied. Thus the displayed relation records the source's qualitative smallness condition, not a specified inequality or limiting procedure.

The source question is: can chemotaxis be the principal mechanism of cell sorting in this model, and what behavior occurs in the regime \(\chi_2\ll 1\)? Here neither ``cell sorting'' nor ``principal mechanism'' is given a mathematical definition. In particular, the record does not specify a criterion for spatial separation, the relevant time scale, or the comparison needed to identify a principal mechanism.

A precise quantified problem remains unresolved at the level of the supplied formulation: the model equations, spatial domain, initial and boundary conditions, solution class, parameter hypotheses, and sorting observable must be specified before a theorem or conjecture can be stated. No particular Keller--Segel system or such hypotheses are imposed in this draft.

The supplied background reports that radial estimates provide regimes of global existence and blow-up, and that blow-up of one species forces blow-up of the other in a cited theorem. The theorem's hypotheses and conclusion are not reproduced, so this implication is retained only as a qualified report of background results, not as an assertion for every two-species Keller--Segel model. The supplied assessment says the small-\(\chi_2\) sorting question is not fully resolved.
\subsection{Short English statement}
The question asks whether chemotaxis can drive cell sorting in a two-species, two-dimensional Keller--Segel model, especially when \(\chi_2\ll 1\). The supplied record does not define the system, the parameter, or a mathematical criterion for sorting, so it does not determine a precise quantified conjecture. Its partially solved label and descriptions of radial results are supplied curation requiring independent review.
\subsection{Sources}
\begin{itemize}
\item Local UnsolvedMath registry, permanent ID \(30003940\), source problem code \texttt{catalogue entry}, title ``Chemotactic Cell Sorting in Keller--Segel Models,'' in the Oberwolfach Reports Open Problems collection. The supplied record gives proposed year 2018 and names no individual proposer.
\item Supplied source citation: \emph{Reactive Flows in Deformable, Complex Media} (2019), Oberwolfach Reports, pp.\ 2385--2473, DOI 10.4171/owr/2018/39. The literature-source label is ``Oberwolfach Report 39/2018.'' \url{https://doi.org/10.4171/owr/2018/39}. Both year designations are retained as supplied.
\item Supplied literature reference: \emph{Sorting Phenomena in a Mathematical Model for Two Species}. \url{https://doi.org/10.48550/arXiv.1704.04179}. The record supplies no authors or theorem identifier.
\item Literature assessment dated exactly \texttt{2026-08-21}: the supplied status is ``partially solved.'' The assessment reports radial global-existence and blow-up regimes, a theorem in which blow-up of one species forces the other, and an incompletely resolved small-\(\chi_2\) sorting question. This is supplied dated curation derived from the OWR record, not independently verified literature assessment. Time-sensitive status and cited mathematical claims require independent verification.
\end{itemize}
\end{document}
